% Section propre à la version complète : gaps entre plateformes (gaps.md, 8e passage)
\subsection{Gaps entre plateformes}
\label{sec:gaps-transverses}

\begin{table}[!htbp]
\centering\footnotesize
\caption{Gaps entre plateformes (repris de \texttt{gaps.md}, 8\textsuperscript{e} passage).}
\label{tab:gaps-transverses}
\begin{tabularx}{\textwidth}{@{}P{4.6cm}P{2.3cm}L@{}}
\toprule
Gap & Statut & Piste \\
\midrule
Pas de synthèse quantitative FPGA / Jetson / MCU à NPU, ni de synthèse sûreté ×
  accélérateur FPGA de détection & étayé (16 surveys) & argument de positionnement pour
  l'état de l'art d'un pré-projet \\
Aucune passerelle MCU--FPGA pour la détection & hypothèse (8 fiches) & porter
  TinyissimoYOLO, ou le LP-YOLO Slim w2a4 de
  \citet{arxiv-2609.24757-npu-accelerator-quantized-real-time}\pp, sur petit FPGA et sur
  MCU, et mesurer l'énergie par inférence à portée égale \\
Aucune comparaison MCU / FPGA / Jetson pour un même algorithme d'apprentissage continu &
  étayé (9 fiches) & porter SLDA ou QLR-CL (étage figé INT8) sur KV260, sur Orin et sur
  GAP9, avec la même mesure de puissance \\
Baseline Jetson Orin mesurée face à une FPGA dans les mêmes conditions & étayé (27 fiches) ;
  partiellement comblé sur AGX Xavier
  \citep{arxiv-2304.11580-framework-benchmarking-real-time-embedded} & même modèle, même
  précision numérique, même méthode de puissance sur KV260 et Orin Nano \\
Anomalie par AE hors FPGA & 2 lignes sur Xavier NX
  \citep{arxiv-2407.17647-energy-efficient-artefact-detection-accelerator}, aucune sur
  MCU & baseline MCU pour
  tout AE d'anomalie porté sur FPGA \\
\bottomrule
\end{tabularx}
\end{table}

La passerelle MCU--FPGA est un sujet à part entière. Côté MCU, les accélérateurs intégrés
(MAX78000, GAP9, STM32N6) atteignent des énergies de l'ordre de la centaine de
microjoules par inférence pour de très petits détecteurs
\citep[§III]{arxiv-2306.00001-tinyissimoyolo-quantized-low-memory-footprint}. Côté FPGA,
les détecteurs compressés du DAC-SDC ont des tailles comparables : 64\,Ko de poids pour
SystemsETHZ \citep[§4.2]{2021-xu-dac-sdc-low-power-object}. Le premier détecteur dataflow
de la base de taille proche d'un MCU est un YOLOv3-tiny allégé en w2a4 par FINN sur XC7Z020, à 2,91\,W au rail 12\,V
\citep[§V-VI]{arxiv-2609.24757-npu-accelerator-quantized-real-time}\pp, mais il n'est pas
porté sur MCU. Aucun travail du corpus ne mesure un même réseau sur les deux. Un MCU à NPU
est pourtant le concurrent le plus direct d'un petit FPGA en détection à très basse consommation \deduit.
